\chapter{Operations \& the web UI}
\label{ch:ops}

\section{Reaching it}
\label{sec:screens}

The bridge serves a web interface on port \code{8080} \emph{inside} the container. It
publishes no host port: the reverse proxy reaches it over a shared Docker network and is the
only thing that can. There is \textbf{no login} --- the trust boundary is the proxy, not this
process.

The interface answers in every state the bridge can be in \emph{once it is running}, and that
is deliberate rather than convenient: since \code{config.toml} became the configuration, a
browser is the only way in, so a server that only ran alongside a healthy session would be
absent from exactly the situations it exists for.

\begin{note}
\textbf{There is no state without a web interface.} A \code{config.toml} that is present and
unusable --- malformed, carrying a key this build does not know, from another schema version,
or unreadable altogether --- does \emph{not} stop the bridge. It comes up, opens no connection
to the broker, emits no birth, publishes nothing, and serves the configuration screen with
every fault rendered against the field it belongs to. \code{/healthz} answers \code{200}, so a
container healthcheck cannot restart away the screen you need.

\textbf{Until 2026-08-06 it refused to start}, and refusing meant no screen --- which made the
commonest first-run mistake, a state directory nobody \code{chown}ed, a restart loop with no
way out that did not involve a shell on the volume. The faults are still printed on
\code{stderr} and in \code{docker compose logs}, each naming the key to edit.

The distinction that matters is \emph{where the repair is}. A file the bridge could read and
rejected is repaired in the form: correct the fields, press \emph{Save}, and the unusable file
is replaced wholesale. A file it could not read \emph{at all} is a filesystem problem --- the
fault says so, and names the \code{chown} --- and no amount of typing in the browser will fix
it, because the bridge cannot write there either.

Correcting the fields and pressing \emph{Save} overwrites the unusable file: a
file this build can diagnose as broken --- malformed TOML, a key it does not know, a
\code{schema\_version} older than its own --- is exactly what the form was rendered from, and
writing it back is the repair.

\textbf{One file is still not overwritten, and deliberately:} one carrying a
\code{schema\_version} \emph{newer} than this build's, which is to say one written by a later
image. It holds settings this build cannot represent, so replacing it would discard them
silently. The refusal names the version, and the answer is to roll the image forward again
rather than to delete the file.
\end{note}

\begin{center}
\begin{tabular}{@{}lll@{}}
\toprule
\textbf{Page} & \textbf{What it is for} & \textbf{Available when} \\
\midrule
\code{/}         & Which state the bridge is in, and why & see below \\
\code{/config}   & Every setting, and the faults in one  & see below \\
\code{/confirm}  & The mapping, to check before publishing & a valid file exists \\
\code{/healthz}  & The machine-facing answer & see below \\
\bottomrule
\end{tabular}
\end{center}

\section{A first run, start to finish}

\begin{enumerate}
  \item Start the container with the credential in \code{.env} and nothing else configured.
        It comes up, publishes nothing, and says so.
  \item Open it in a browser and go to \textbf{Configuration}. Fill in the Sparkplug group
        and node, the broker, the publish period, and one meter --- its name, its smart-me
        device id and its serial. Tick \emph{Published}.
  \item Save. If anything is wrong, \textbf{every fault is shown at once}, each beside the
        field it belongs to, in the same words the bridge would have used had it read the
        value from the file. There is one set of rules, not two: a value this form accepts
        is a value the bridge will start on.
  \item Go to \textbf{Confirm the meter mapping}. Check each \textbf{serial against its
        topic}. This is the one error no machine can catch --- a mapping can be
        well-formed, unique, complete and pointed at the wrong meter --- so it is a separate
        page and a separate click, never a checkbox on the save form.
  \item Confirm. The bridge starts publishing. \textbf{You do not restart anything.}
\end{enumerate}

\begin{note}
There is \textbf{no credential field} on any screen, and that is not an oversight: the
smart-me credential lives in the environment and never descends to disk, so the strongest
form of \emph{never rendered} turned out to be \emph{never present}. If a fault names
\code{SMARTME\_CLIENT\_ID} or \code{SMARTME\_CLIENT\_SECRET} it is reported as an
environment problem, not drawn as a box you can type in --- because you cannot fix it here.
\end{note}

\begin{warning}
\textbf{The state directory must be writable by the container's user.} It runs as uid
\code{10002}, so a bind mount created by anyone else is read-only to it and the save will
fail. The screen says so plainly --- \emph{``the configuration was NOT written''} --- and
changes nothing, rather than reporting a success that never reached the disk. The fix is
\code{chown 10002:10002} on the directory, once, before the first run.
\end{warning}

\section{Changing a setting on a running bridge}

Not every setting costs the same, and the screen says which before you believe it happened.
After a save it reports one of four verdicts:

\begin{center}
\begin{tabular}{@{}ll@{}}
\toprule
\textbf{Verdict} & \textbf{What actually happened} \\
\midrule
Saved, and in force now          & The running bridge is already using it \\
One device re-announced          & Same session, same \code{bdSeq}, one fresh DBIRTH \\
Saved --- but NOT in force       & Needs a new Sparkplug session; restart to apply \\
Saved --- but NOT in force       & Needs a process restart; the setting is on disk \\
\bottomrule
\end{tabular}
\end{center}

The distinction is the point. Changing the broker, the group id or the node id is a
\emph{new session} by definition, and this build cannot open one without being restarted.
Changing the web UI port likewise: a listener cannot move without dropping what is connected
to it. The screen names each setting still waiting, and re-reads what is genuinely in force
rather than echoing back what you typed --- so it never reports a change the bridge did not
make.

\begin{warning}
\textbf{Withdrawing a confirmation still costs a restart.} Editing the mapping --- the meters,
or \code{group\_id} or \code{node\_id}, which are in every topic --- clears
\code{mapping\_confirmed}, correctly, since it is a different mapping and nobody has
looked at it. But a bridge that is already publishing keeps its session until it is
restarted. Confirming moves you forward without a restart; going back does not.
\end{warning}

\section{Submissions from elsewhere are refused}

Because there is no login, any page open in your browser could otherwise post to the bridge
if that browser can reach it --- a reconfiguration needs no credential to steal when there
are none. Every submission is therefore checked against the page it claims to come from, and
one that came from somewhere else is refused with \emph{Refused}, changing nothing.

A caller that sends no \code{Origin} at all --- \code{curl}, a script, the container smoke
test --- is allowed: it is not a browser, and refusing it would break the headless path
while stopping no attack. See
\code{docs/adr/0024-the-config-ui-refuses-submissions-from-other-origins.md}.

\section{What is not here yet}

\stub{per-meter live values, dual source/sink health, the ``state of the bridge''
orientation screen, culprit classification and on-demand end-to-end validation are later
stories in this epic. Today the interface configures the bridge and reports which state it
is in; it does not yet show what the meters are reading.}

\section{The health endpoint}

\code{GET /healthz} is the machine-facing answer and is described with the rest of
diagnosis in chapter~\ref{ch:trouble}. The short version: the \emph{status code} answers
``is this process worth keeping?'' and the \emph{body} answers ``is it doing anything?''.
They are two different questions, and a bridge that is deliberately silent --- unconfigured,
or waiting for you to confirm --- answers \code{200}, because the healthcheck restarts the
container and restarting one would destroy the very screen needed to configure it.
